\chapter{Objectives}\label{chap:objectives}

This chapter presents the objectives that guide the development of this dissertation. They are organized into a general objective, which defines the overarching research goal, and a set of specific objectives that outline the individual steps required to achieve it.

\section{General Objective}

The primary goal of this research is to develop and evaluate machine learning models capable of automating the detection of Equatorial Plasma Bubbles in airglow images at the OI 630 nm emission from the EMBRACE INPE program, contributing to the improvement of operational space weather monitoring in Brazil.

\section{Specific Objectives}

To achieve this goal, the following specific objectives guide the development of this dissertation:

\begin{enumerate}
    \item Conduct a systematic literature review on the application of machine learning techniques for detecting ionospheric irregularities and structures in airglow imagery, identifying the architectures, preprocessing strategies, and evaluation metrics currently in use.
    \item Build a labeled database of all-sky airglow images from the EMBRACE INPE archives, containing examples of both EPB presence and absence across different observational conditions, seasons, and levels of solar activity.
    \item Design and implement a physics-informed preprocessing pipeline that models and removes background illumination, enhances local contrast through CLAHE, and normalizes pixel intensities to isolate the physical signatures of plasma depletions.
    \item Extract spatial features from the preprocessed images using Two-Dimensional Principal Component Analysis (2DPCA), preserving the two-dimensional correlations while substantially reducing data dimensionality.
    \item Train and evaluate a stacking ensemble classifier combining Support Vector, XGBoost, and Random Forest algorithms, integrated with SMOTE to address class imbalance, for the automatic classification of EPB events.
    \item Assess the interpretability of the trained model through SHapley Additive exPlanations (SHAP), verifying that the classifier decisions are driven by the physical morphology of plasma depletions rather than environmental artifacts.
    \item Validate the temporal stability of the system over continuous nocturnal observation sequences, ensuring reliable performance under real-time monitoring conditions.
    \item Compare the proposed architecture against alternative machine learning approaches using quantitative metrics such as precision, recall, F1 score, and Intersection over Union, complemented by domain expert validation.
\end{enumerate}